\appendix
\section*{附录：数据来源与复现}
\addcontentsline{toc}{section}{附录：数据来源与复现}
\noindent\textbf{数据来源与复现。}
文中数值与图表来自本文定义的抽象随机实验，
由精确组合计算和Python实体排列模拟生成。
107组主实验、基准复验及精度实验共使用$30\,260\,000$次排列。
复现环境为Python 3.13.11、NumPy 2.5.4、Matplotlib 3.11.2。
项目保存配置、整数直方图、理论分数、种子和源文件摘要；
正文表格由这些既有结果自动提取，7组矢量图沿用已核验数据与绘图程序，
并为A4印刷放大标注字号；原始图另行保留。
论文目录中的复现说明列出了完整命令及数据对应关系。

\medskip
\noindent\textbf{推广计数。}
对多副本A/B/C构筑，若设置阶段相应计数为$i,j,\ell$，
S计数为$u$，则成功要求$i,j,\ell\ge1$，剩余可消耗数为
\[
M=(a-i)+(b-j)+(c-\ell)+(s-u).
\]
设置计数的组合权重为
\[
q(i,j,\ell,u)=
\frac{\binom ai\binom bj\binom c\ell\binom su
\binom{d+t}{n-i-j-\ell-u}}{\binom Nn}.
\]
给定计数后，收益分布为$\operatorname{Hypergeom}(N-n,M,h)$；
按权重混合即推广式\eqref{eq:joint}，并保留单对A/B激活的固定收益系数8。
